\section{A follow-on theorem enabled by the formal development}
\label{sec:example-transport}

The source correction asks whether direct rotation minimizes the cost
of transporting a subspace.  This question differs from the perturbation
estimates: here the subspaces are fixed and the cost is measured over
all orthogonal maps between them.  The formalization supplied the
principal-angle, real/complex, and unitarily invariant norm interfaces
needed to examine this larger family of admissible transports.

Let $K,L\subseteq\mathbb R^n$ have the same dimension and let
$\theta_1,\ldots,\theta_m\in(0,\pi/2]$ be their positive principal angles,
counted with multiplicity.  A matching $M$ on the angle labels consists
of disjoint unordered pairs; labels outside the pairs are unmatched.
Define the singleton and paired costs
\begin{equation}\label{eq:nuclear-matching-costs}
 s(\theta)=\min\{4\sin(\theta/2),2\},\qquad
 e(\theta,\psi)=4\sin((\theta+\psi)/2).
\end{equation}

\begin{theorem}[Verified real nuclear matching reduction]
\label{thm:real-matching-example}
For equal-dimensional real subspaces $K,L$ as above,
\begin{equation}\label{eq:real-nuclear-match}
\min_{\substack{W\in O(n)\\W(K)=L}}\|I-W\|_1
=
\min_{M\text{ a matching}}
\left(
\sum_{\{i,j\}\in M}e(\theta_i,\theta_j)
+
\sum_{i\text{ unmatched}}s(\theta_i)
\right).
\end{equation}
Both minima are attained.  The equality includes repeated angles,
right angles, and common ambient directions; the left side has no
principal-plane or determinant constraint.
\end{theorem}

The Lean development proves the unrestricted equality in
\cref{eq:real-nuclear-match}, including comparison with arbitrary
ambient orthogonal transporters.  Pairing two real principal planes
allows a displacement pattern unavailable to a single real plane;
this is the mechanism visible in the four-dimensional counterexample.
For $\theta_1=\theta_2=\pi/6$, the canonical direct rotation costs
$8\sin(\pi/12)>2$, whereas the paired competitor costs
$4\sin(\pi/6)=2$.

The complex-phase construction behind this example was already present
in Davis and Kahan \citep[Example~4.2]{davis1970rotation}.  The general
unitary minimal-distortion problem was studied by Bolotnikov, Li, and
Rodman \citep{bolotnikov2001minimal}.  Theorem~\ref{thm:real-matching-example} gives an exact global
weighted-matching characterization for arbitrary finite-dimensional
real orthogonal transporters; the earlier complex example and the
minimal-distortion results provide its geometric and optimization
context.  The conventional proof and additional Schatten and
smoothed-penalty results are developed in the separate
spectral-transport manuscript.  The matching theorem is presented
here as a mathematical result enabled by the formalization of the
Davis--Kahan foundations.
